% Anexo de la sección C. Requiere en el preámbulo: array, booktabs, longtable.

\section*{Anexo C. Modelo dimensional: matrices y diccionario de datos}

\subsection*{C.1 Matriz de hechos y calificadores}

Incluye las cuatro preguntas del README y las dos que la contraparte planteó en la reunión del 4 de septiembre. En cada una, el verbo de medición es el hecho y cada ``por'' es un calificador.

\begin{center}
\small
\begin{tabular}{@{}>{\raggedright\arraybackslash}p{0.40\textwidth}>{\raggedright\arraybackslash}p{0.25\textwidth}>{\raggedright\arraybackslash}p{0.28\textwidth}@{}}
\toprule
Pregunta del negocio & Hecho (qué se mide) & Calificadores (por dónde se corta) \\
\midrule
¿El precio por litro subió por precio o por mezcla? (README 2, la elegida) & Monto neto, litros & Tier, semestre, mercado \\
¿Las exportaciones crecen en pesos? (README 1) & Monto neto en dólares y pesos, tipo de cambio & Mes, semestre, mercado \\
¿Qué mercados crecieron en volumen y con qué margen? (README 3) & Litros, monto neto, costo total & Mercado, semestre, tier \\
¿Qué fecha usar para convertir y qué ingreso comparar con el costo? (README 4) & Monto neto, costo total, tipo de cambio & Fecha de factura, fecha de embarque \\
¿Cuántos litros vendemos de cada cepa y de cada tier? (reunión) & Litros & Cepa, tier, mercado, mes \\
¿En qué mercados poner más esfuerzo? (reunión) & Litros, monto neto & Mercado, cliente, tier \\
\bottomrule
\end{tabular}
\end{center}

Lo que revela la matriz: el monto neto y los litros se repiten en casi todas las preguntas, así que comparten grano y van en la misma tabla de hechos. Tiempo, mercado y producto (tier y cepa) aparecen en varias preguntas: son las dimensiones conformadas. La fecha de embarque pertenece a otro proceso, el embarque, y por eso no entra a esta estrella.

\subsection*{C.2 Matriz de bus}

\begin{center}
\small
\begin{tabular}{@{}lccccc@{}}
\toprule
Proceso & Tiempo & Producto & Cliente & Mercado & Incoterm \\
\midrule
Venta (modelada en la Entrega 1) & $\bullet$ & $\bullet$ & $\bullet$ & $\bullet$ & \\
Embarque (candidato, JSON) & $\bullet$ & $\bullet$ & & $\bullet$ & $\bullet$ \\
Cobranza (candidato, sin fechas de pago) & $\bullet$ & & $\bullet$ & $\bullet$ & \\
\bottomrule
\end{tabular}
\end{center}

Tiempo y mercado se repiten en los tres procesos: si cada área los definiera distinto, ningún cruce entre venta, embarque y cobranza cuadraría. Por eso se definen una sola vez.

\subsection*{C.3 Dimensiones que cambian}

Las fuentes son una foto sin historia de cambios, así que todas las dimensiones se cargan con tipo 1 (sobrescribir). Si en adelante el maestro de productos cambiara el costo o los litros por caja de un producto, habría que versionar ese atributo (tipo 2), para que las ventas pasadas conserven el valor vigente al momento de la venta.

\subsection*{C.4 Diccionario de datos}

% Generado por src/generar_diccionario.py desde src/esquema.py. No editar a mano.

\subsubsection*{\texttt{fact\_venta}}
Hecho transaccional. Una fila por línea de venta facturada.

{\footnotesize
\begin{longtable}{@{}>{\raggedright\arraybackslash}p{0.19\textwidth}>{\raggedright\arraybackslash}p{0.1\textwidth}>{\raggedright\arraybackslash}p{0.085\textwidth}>{\raggedright\arraybackslash}p{0.205\textwidth}>{\raggedright\arraybackslash}p{0.27\textwidth}@{}}
\toprule
Columna & Rol & Aditiv. & Descripción & Origen y regla \\
\midrule
\endhead
\texttt{venta\_id} & degenerada & — & Número de la línea en el ERP. & ventas.\allowbreak{}venta\_id; deduplicado tras normalizar formatos \\
\texttt{fecha\_key} & llave & — & Fecha de factura. & ventas.\allowbreak{}fecha, leída en sus 4 formatos \\
\texttt{producto\_anada\_key} & llave & — & Producto y añada vendidos. & ventas.\allowbreak{}producto\_id + ventas.\allowbreak{}anada \\
\texttt{cliente\_key} & llave & — & Cliente que compra. & ventas.\allowbreak{}cliente\_id \\
\texttt{mercado\_key} & llave & — & Mercado de destino. & ventas.\allowbreak{}mercado\_id (coincide con clientes.\allowbreak{}mercado\_id) \\
\texttt{cajas} & medida & aditiva & Cajas vendidas; negativas si la línea es merma. & ventas.\allowbreak{}cajas; si vienen en botellas (12 o 10 veces las cajas del embarque), se dividen por el factor \\
\texttt{litros} & medida & aditiva & Litros vendidos. & cajas $\times$ dim\_producto.\allowbreak{}litros\_por\_caja \\
\texttt{monto\_bruto\_clp} & medida & aditiva & Monto antes de descuento, en pesos. & cajas $\times$ precio unitario; si la moneda es USD, $\times$ tipo\_cambio \\
\texttt{monto\_bruto\_usd} & medida & aditiva & Monto antes de descuento, en dólares. & cajas $\times$ precio unitario; si la moneda es CLP, $\div$ tipo\_cambio \\
\texttt{monto\_neto\_clp} & medida & aditiva & Ingreso neto de descuento, en pesos. & cajas $\times$ precio unitario $\times$ (1 $-$ descuento), convertido con tipo\_cambio \\
\texttt{monto\_neto\_usd} & medida & aditiva & Ingreso neto de descuento, en dólares. & cajas $\times$ precio unitario $\times$ (1 $-$ descuento), convertido con tipo\_cambio \\
\texttt{costo\_total\_clp} & medida & aditiva & Costo de producción de lo vendido, en pesos. & \textbar{}cajas\textbar{} $\times$ dim\_producto.\allowbreak{}costo\_caja\_clp; en la merma el costo no se reversa \\
\texttt{tipo\_cambio} & contexto & no aditiva & Dólar observado usado para convertir la línea. & API mindicador.\allowbreak{}cl, fecha de factura; fines de semana y feriados con el último día hábil previo \\
\texttt{es\_merma} & indicador & aditiva & 1 si la línea es merma (cajas negativas). & cajas < 0; merma según la contraparte (4 sep) \\
\texttt{monto\_recalculado} & indicador & aditiva & 1 si el monto registrado no cuadraba y se recalculó. & \textbar{}monto registrado $-$ calculado\textbar{} > 0,5 \% \\
\texttt{cajas\_corregidas} & indicador & aditiva & 1 si las cajas venían en botellas. & cruce con origen\_embarques.\allowbreak{}json \\
\texttt{moneda\_imputada} & indicador & aditiva & 1 si la moneda venía vacía. & ventas.\allowbreak{}moneda\_id nulo; se imputa por mercado \\
\bottomrule
\end{longtable}
}

\subsubsection*{\texttt{dim\_tiempo}}
Una fila por día calendario de 2025.

{\footnotesize
\begin{longtable}{@{}>{\raggedright\arraybackslash}p{0.19\textwidth}>{\raggedright\arraybackslash}p{0.1\textwidth}>{\raggedright\arraybackslash}p{0.085\textwidth}>{\raggedright\arraybackslash}p{0.205\textwidth}>{\raggedright\arraybackslash}p{0.27\textwidth}@{}}
\toprule
Columna & Rol & Aditiv. & Descripción & Origen y regla \\
\midrule
\endhead
\texttt{fecha\_key} & llave & — & Llave AAAAMMDD. & generada \\
\texttt{fecha} & atributo & — & Fecha. & calendario \\
\texttt{dia\_semana} & atributo & — & Lunes a domingo. & calendario \\
\texttt{semana} & atributo & — & Semana del año. & calendario \\
\texttt{mes} & atributo & — & Mes (número y nombre). & calendario \\
\texttt{trimestre} & atributo & — & Trimestre. & calendario \\
\texttt{semestre} & atributo & — & Primer o segundo semestre. & calendario \\
\texttt{anio} & atributo & — & Año. & calendario \\
\texttt{es\_dia\_habil} & atributo & — & 1 si es día hábil. & 1 si la API del dólar publica valor ese día \\
\texttt{es\_feriado} & atributo & — & 1 si es feriado. & día de semana sin valor publicado en la API \\
\bottomrule
\end{longtable}
}

\subsubsection*{\texttt{dim\_producto}}
Una fila por producto y añada (24 productos $\times$ 4 añadas = 96 filas).

{\footnotesize
\begin{longtable}{@{}>{\raggedright\arraybackslash}p{0.19\textwidth}>{\raggedright\arraybackslash}p{0.1\textwidth}>{\raggedright\arraybackslash}p{0.085\textwidth}>{\raggedright\arraybackslash}p{0.205\textwidth}>{\raggedright\arraybackslash}p{0.27\textwidth}@{}}
\toprule
Columna & Rol & Aditiv. & Descripción & Origen y regla \\
\midrule
\endhead
\texttt{producto\_anada\_key} & llave & — & Llave producto-añada. & generada \\
\texttt{producto} & atributo & — & Nombre del vino. & origen\_costos\_productos.\allowbreak{}csv: producto \\
\texttt{cepa} & atributo & — & Cepa, con el nombre oficial del catálogo del ERP. & CSV: cepa, normalizada desde sus alias y conformada a origen\_vina.\allowbreak{}sqlite: cepas.\allowbreak{}cepa; desempate en la bitácora \\
\texttt{tipo\_vino} & atributo & — & Tinto o blanco. & origen\_vina.\allowbreak{}sqlite: cepas.\allowbreak{}tipo \\
\texttt{tier} & atributo & — & Varietal, Reserva o Gran Reserva. & CSV: tier \\
\texttt{formato} & atributo & — & 12 $\times$ 750 ml o 6 $\times$ 1500 ml. & CSV: formato \\
\texttt{litros\_por\_caja} & atributo & — & Litros que rinde una caja: la definición de caja a litro. & CSV: litros\_por\_caja (9,0 para ambos formatos) \\
\texttt{anada} & atributo & — & Año de cosecha. & CSV: anada \\
\texttt{costo\_caja\_clp} & atributo & — & Costo unitario de producción por caja. & CSV: costo\_caja\_clp; si falta, regla exacta (tier, añada) \\
\texttt{costo\_imputado} & atributo & — & 1 si el costo se imputó por regla. & generado \\
\texttt{activo} & atributo & — & Si el producto sigue en catálogo. & CSV: activo (si/SI/Sí/1 = verdadero) \\
\bottomrule
\end{longtable}
}

\subsubsection*{\texttt{dim\_cliente}}
Una fila por cliente (distribuidor o importador).

{\footnotesize
\begin{longtable}{@{}>{\raggedright\arraybackslash}p{0.19\textwidth}>{\raggedright\arraybackslash}p{0.1\textwidth}>{\raggedright\arraybackslash}p{0.085\textwidth}>{\raggedright\arraybackslash}p{0.205\textwidth}>{\raggedright\arraybackslash}p{0.27\textwidth}@{}}
\toprule
Columna & Rol & Aditiv. & Descripción & Origen y regla \\
\midrule
\endhead
\texttt{cliente\_key} & llave & — & Llave del cliente. & generada \\
\texttt{cliente\_id} & atributo & — & Código del ERP. & origen\_vina.\allowbreak{}sqlite: clientes.\allowbreak{}cliente\_id \\
\texttt{nombre} & atributo & — & Razón social. & clientes.\allowbreak{}nombre \\
\texttt{tipo} & atributo & — & Distribuidor o importador. & clientes.\allowbreak{}tipo \\
\bottomrule
\end{longtable}
}

\subsubsection*{\texttt{dim\_mercado}}
Una fila por mercado de destino. País y zona son niveles de una misma jerarquía.

{\footnotesize
\begin{longtable}{@{}>{\raggedright\arraybackslash}p{0.19\textwidth}>{\raggedright\arraybackslash}p{0.1\textwidth}>{\raggedright\arraybackslash}p{0.085\textwidth}>{\raggedright\arraybackslash}p{0.205\textwidth}>{\raggedright\arraybackslash}p{0.27\textwidth}@{}}
\toprule
Columna & Rol & Aditiv. & Descripción & Origen y regla \\
\midrule
\endhead
\texttt{mercado\_key} & llave & — & Llave del mercado. & generada \\
\texttt{pais} & atributo & — & País de destino. & origen\_vina.\allowbreak{}sqlite: mercados.\allowbreak{}pais \\
\texttt{zona} & atributo & — & Zona comercial. & mercados.\allowbreak{}zona\_id $\rightarrow$ zonas.\allowbreak{}zona; zona 8 inexistente (Corea del Sur) se asigna a Asia \\
\texttt{canal} & atributo & — & Nacional o Exportación. & mercados.\allowbreak{}canal \\
\texttt{hemisferio} & atributo & — & Norte o sur. & derivado del país (Chile y Brasil: sur) \\
\texttt{moneda\_facturacion} & atributo & — & CLP o USD. & monedas.\allowbreak{}moneda; en los datos la moneda queda determinada por el mercado (Chile CLP, resto USD); el ETL verifica que no haya excepciones \\
\bottomrule
\end{longtable}
}

